\subsection{PART A – 7Cs of Communication}

\indent 

Rewrite the following sentence into a better one using the 7Cs of Communication: “Hey, project’s late, not sure when it’ll be done. Will let you know”

Possible good version: “Hi Alex, the login module is delayed due to integration issues with the API. We need two more days to resolve and test. I’ll provide an update by Thursday 5 PM and escalate to the security team if needed.”

\subsection{PART B – Stakeholder Mapping}

\indent 

\textbf{Case Scenario}: A university is launching a new online learning platform

\noindent\textbf{In groups, identify at least 5 stakeholders for this project. Place them into a Power–Interest Matrix}

High Power, High Interest (Manage Closely): University administrators, government regulators.
High Power, Low Interest (Keep Satisfied): IT directors, funding bodies.
Low Power, High Interest (Keep Informed): Students, lecturers.
Low Power, Low Interest (Monitor): Technical support staff not directly on the project.

\subsection{PART C – Case Study Discussion}

\indent 

An IT team is developing a mobile app for a healthcare provider. Developers and testers are clashing, doctors complain they’re not consulted, and project updates are unclear to management.

\bigskip

\noindent\textbf{Where is communication breaking down?}

Unclear project updates to management; no feedback channel for doctors.

\bigskip

\noindent\textbf{What teamwork challenges can you spot?}


Developers and testers in conflict; poor documentation sharing.


\bigskip

\noindent\textbf{How would you manage stakeholders here?}


Executives (sponsors):
\begin{itemize}
    \item High power, high interest → Manage closely.
    \item Provide regular progress updates (e.g., fortnightly demos, clear timelines using SBAR).
    \item Address their need for faster release by explaining the trade-offs with compliance/security.
\end{itemize}

\bigskip 

Doctors (end-users):
\begin{itemize}
    \item High interest, medium power → Keep informed and engaged.
    \item Involve them in requirements workshops and usability testing so the final app supports real workflows.
    \item Provide clear channels (surveys, feedback sessions) for their input.
\end{itemize}

Compliance officers (regulators):
\begin{itemize}
    \item High power, high interest → Manage closely.
    \item Keep them updated with compliance documentation and involve them early in testing to avoid last-minute blockers.
\end{itemize}

Testers \& Developers (internal team):
\begin{itemize}
    \item Medium power, high interest → Empower through teamwork.
    \item Improve collaboration via shared documentation, standups, retrospectives.
\end{itemize}

\bigskip

\noindent\textbf{How do communication, teamwork, and stakeholder management come together to fix this?}

Communication provides clarity: Using structured updates (7 Cs or SBAR), the team keeps management and compliance informed while giving doctors a voice in development. This prevents misunderstandings and builds trust.

Teamwork provides collaboration: Developers and testers align through Agile ceremonies (daily stand-ups, retrospectives, shared tools), resolving conflicts and boosting productivity.

Stakeholder management provides direction: By mapping stakeholders in the Power–Interest Matrix, the team prioritizes who to engage closely (executives, compliance) and who to keep informed (doctors).


